\documentclass[11pt]{article}
\usepackage[margin=1in]{geometry}
\usepackage[T1]{fontenc}
\usepackage{lmodern,microtype,hyperref,enumitem,amsmath,amssymb,xcolor,booktabs}
\usepackage[numbers,sort&compress]{natbib}
\hypersetup{colorlinks=true,urlcolor=blue,citecolor=blue,linkcolor=blue}
\title{Q07: Comprehension Check}
\author{DLS 2026 --- Kukla}
\date{October 2026}
\begin{document}
\maketitle
\textbf{Comprehension check.} A lab has a fixed compute budget and plans to spend it all on a very large
model trained on a relatively small dataset. Citing the compute-optimal result, what is likely wrong with
this plan, and what does $C \approx 6ND$ tell them to do instead?

\vspace{0.8cm}\hrule\vspace{4pt}
{\small\textbf{Answer.} The plan over-parameterizes and under-trains: Chinchilla showed that for a fixed
compute budget, loss is minimized by scaling model size $N$ and training tokens $D$ \emph{together}
(roughly equally), not by maximizing $N$ alone. Since $C \approx 6ND$ is fixed, pouring budget into $N$
forces $D$ down, leaving the big model undertrained and achieving \emph{higher} loss than a smaller model
trained on more tokens at the same cost. They should rebalance toward a smaller $N$ and a larger $D$ to
hit the compute-optimal point.}
\end{document}
